%% 01_introduction.tex
\section{Introduction}
\label{sec:intro}

Experimental asset markets in the tradition of
\citet{smith1988bubbles} have produced one of the most replicated
departures from rational-expectations pricing in all of economics:
even when subjects know, in closed form, that the fundamental value
of a finite-lived asset declines deterministically with the
remaining dividend stream, transaction prices systematically
overshoot the fundamental before collapsing near the terminal
period. \citet{dufwenberg2005bubbles}---hereafter
DLM~2005---established a by now canonical result in this literature:
when two-thirds or more of the subject pool consists of traders who
have played the game before, the bubble~disappears. Experience, rather
than cognition or incentives, is what eventually kills the bubble.

A parallel and much more recent literature replaces the human
subjects with large language models. \citet{lopezlira2025llm}
proposes an expected-utility framework in which each LLM-flavoured
agent carries a persistent valuation bias $b_i$, a risk preference
over wealth $U\in\{\Uloving,\Uneutral,\Uaverse\}$, and a belief-blending
channel that lets agents exchange private valuations and update a
trust score for each sender by exponential moving average. Lopez-Lira's
agents optimise expected utility over a small action set---hold,
cross the bid, cross the ask, post passively---and the resulting
population can be dialled between skeptical and naive, truthful and
deceptive, risk-averse and risk-loving by simple categorical knobs.
The Lopez-Lira design is a natural vehicle for studying~how a
heterogeneous LLM population prices a DLM asset, but it leaves open a
psychological question the original DLM protocol never asked: if one
replaces the deterministic prior~$\FV_t\cdot(1+b_i+\varepsilon)$ with
a private subjective valuation drawn from a language model that is
told who the agent is and what the market looks like, does the asset
then migrate, over the course of the run, into the hands of the
agent with the highest such valuation?

This paper contributes a computational testbed in which the two
designs share one continuous double-auction market, and a new
paradigm---which I label \emph{Wang~2026}---that seeds each utility
agent's prior with exactly one ChatGPT call at run~start. Concretely,
when the user selects Wang~2026 and supplies an API key, the
simulator fires one chat completion per utility agent in parallel,
using a prompt that encodes the agent's risk preference, bias mode,
and the presence of DLM experienced counterparties in the same
market. The returned number, clamped to
$[0.25\,\FV_0,\,1.75\,\FV_0]$, is written to the agent's
\texttt{psychAnchor} slot and consumed on the first decision tick in
place of the deterministic prior; subsequent ticks revert to the
Lopez-Lira prior plus belief~blending. In the notation of the
replicated model,
\begin{equation}
\hat{V}_i \;\leftarrow\; \mathrm{ChatGPT}(\text{agent}_i),
\qquad
\mathrm{prior}_{t=0}^{(i)} \;=\; \hat{V}_i\cdot(1+\varepsilon_i),
\qquad \hat{V}_i\in[0.25\,\FV_0,\,1.75\,\FV_0].
\label{eq:anchor}
\end{equation}
The deterministic fallback---used when no key is supplied or the
network call fails---drops equation~\eqref{eq:anchor} and resumes the
Lopez-Lira prior $\FV_t\cdot(1+b_i+\varepsilon)$ unchanged. Because
the run then proceeds through a deterministic seeded PRNG, every
simulation can be replayed identically with or without the anchor
channel.

The research question is easy to instrument. Let
$i^\star=\arg\max_{i\in U}\hat{V}_i$ be the utility agent with the
highest anchored valuation and
$j^\star(t)=\arg\max_i\text{inventory}_i(t)$ the top holder at
tick~$t$; a \emph{match} at~$t$ is
$\mathbb{1}\{j^\star(t)=i^\star\}$ and the \emph{normalised gap} is
$g(t)=(\hat V_{i^\star}-\hat V_{j^\star(t)})/\hat V_{i^\star}$.
Both quantities are reported every tick alongside Haessel~$R^2$,
normalised deviation, amplitude, turnover, and allocative
efficiency, so classical bubble statistics stay visible beside the
new Wang metric block.

Two features distinguish the Wang~2026 paradigm from a na\"ive
``LLM-as-trader'' pipeline. First, the language model is consulted
exactly once at run~start; after the anchor is absorbed, the agent is
again a pure expected-utility maximiser and the mulberry32 RNG alone
controls the downstream dynamics. Second, the anchor is a
\emph{valuation} and not an order: the agent's quoting rule, fill
probability, and utility function are unchanged, and only its starting
prior is moved. These two choices isolate the specific question
``does psychology predict ownership'' from the adjacent question
``does an LLM trade~well.''

Section~\ref{sec:lit} places the simulator in the literature on
experimental asset markets and LLM-based agents.
Section~\ref{sec:design} documents the market structure and the
sampling stage. Section~\ref{sec:agents} gives the precise decision
rule for each agent type. Section~\ref{sec:innovation} presents the
Wang~2026 extension: the anchor call, the clamp, the consumption
path, and the research-question instrumentation.
Section~\ref{sec:results} reports preliminary findings, and
Section~\ref{sec:conclusion} concludes.
